% Part: incompleteness
% Chapter: representability-in-q
% Section: minimization-representable

\documentclass[../../../include/open-logic-section]{subfiles}

\begin{document}

\olfileid{inc}{req}{min}
\olsection{کمینه‌سازی منظم در~$\Th{Q}$ بازنمایی‌پذیر است}

جست‌وجوی نامحدود را در نظر بگیریم. فرض کنید $g(x, z)$ منظم و در
$\Th{Q}$ بازنمایی‌پذیر باشد، مثلاً با !!{formula}~$!A_g(x, z, y)$.
بگذارید $f$ با $f(z) = \umin{x}{[g(x, z) = 0]}$ تعریف شود. می‌خواهیم
!!a{formula}ی چون~$!A_f(z, y)$ بیابیم که~$f$ را بازنمایی کند. مقدار~$f(z)$
آن عدد~$x$ است که (الف) در $g(x, z) = 0$ صدق می‌کند و (ب) کوچک‌ترین
چنین عددی است؛ یعنی برای هر $w < x$، $g(w, z) \neq 0$. پس انتخاب زیر
طبیعی است:
\[
!A_f(z,y) \ident !A_g(y, z, \Obj 0) \land \lforall[w][(w < y \lif \lnot
  !A_g(w, z, \Obj 0))].
\]
البته در حالت کلی باید $z$ را با $z_0$، \dots، $z_k$ جایگزین کنیم.

این اثبات نیز چند لم دربارهٔ اموری در بر خواهد داشت که~$\Th{Q}$ برای
اثباتشان به‌اندازهٔ کافی نیرومند است.

\begin{lem}
\ollabel{lem:succ} برای هر !!{constant}~$a$ و هر عدد طبیعی~$n$،
\[
\Th{Q} \Proves \eq[(a' + \num n)][(a + \num n)'].
\]
\end{lem}

\begin{proof}
اثبات، طبق معمول، با استقرا روی~$n$ است. در حالت پایه، $n = 0$، باید
نشان دهیم $\Th{Q}$ فرمول $\eq[(a' + \Obj 0)][(a + \Obj 0)']$ را ثابت
می‌کند. اما داریم:
\begin{align}
  \Th{Q} & \Proves \eq[(a' + \Obj 0)][a'] \quad \text{بنا بر اصل موضوعهٔ $Q_4$}
  \ollabel{step1}\\
  \Th{Q} & \Proves  \eq[(a + \Obj 0)][a] \quad \text{بنا بر اصل موضوعهٔ $Q_4$}
  \ollabel{step2} \\
  \Th{Q} & \Proves \eq[(a + \Obj 0)'][a'] \quad \text{بنا بر \olref{step2}}
  \ollabel{step3} \\
  \Th{Q} & \Proves \eq[(a' + \Obj 0)][(a + \Obj 0)'] \quad
  \text{بنا بر \olref{step1} و \olref{step3}.}\notag
\end{align}
در گام استقرا می‌توانیم فرض کنیم نشان داده‌ایم $\Th{Q}
\Proves \eq[(a' + \num n)][(a + \num n)']$. چون
$\num{n+1}$ همان~$\num{n}'$ است، باید نشان دهیم $\Th{Q}$ فرمول
$\eq[(a' + \num{n}')][(a + \num{n}')']$ را ثابت می‌کند. داریم:
\begin{align}
  \Th{Q} & \Proves \eq[(a' + \num n')][(a' + \num n)'] \quad
  \text{بنا بر اصل موضوعهٔ $!Q_5$} \ollabel{step5}\\
  \Th{Q} & \Proves \eq[(a' + \num n')][(a + \num n')'] \quad
  \text{بنا بر فرض استقرا، اصل موضوعهٔ $!Q_5$ و منطق برابری} \ollabel{step6}\\
  \Th{Q} & \Proves \eq[(a' + \num n)'][(a + \num n')'] \quad
  \text{بنا بر \olref{step5} و \olref{step6}.} \notag
\end{align}
\end{proof}

باز هم شایان ذکر است که این ادعا ضعیف‌تر از آن است که بگوییم~$\Th{Q}$
$\lforall[x][\lforall[y][(x' + y) = (x + y)']]$ را ثابت می‌کند.
با آنکه این !!{sentence} در~$\Struct{N}$ صادق است، $\Th{Q}$ آن را ثابت
نمی‌کند.

\begin{lem}
\ollabel{lem:less-zero}
$\Th{Q} \Proves \lforall[x][\lnot x < \Obj 0]$.
\end{lem}

\begin{proof}
اثبات را به‌طور غیرصوری ارائه می‌کنیم (یعنی فقط راهنمایی‌هایی دربارهٔ
چگونگی ساختن !!{derivation} صوری می‌دهیم).

باید $\lnot a < \Obj 0$ را برای~$a$ دلخواه ثابت کنیم. بنا بر تعریف
$<$، باید $\lnot \lexists[y][\eq[(y'+a)][\Obj 0]]$ را در~$\Th{Q}$
ثابت کنیم. $\lexists[y][\eq[(y'+a)][\Obj 0]]$ را فرض می‌کنیم و
تناقضی به دست می‌آوریم. فرض کنید $\eq[(b'+a)][\Obj 0]$. با استفاده
از~$!Q_3$ داریم $\eq[a][\Obj 0] \lor \lexists[y][\eq[a][y']]$.
دو حالت را از هم متمایز می‌کنیم.

حالت ۱: $\eq[a][\Obj 0]$ برقرار است. از $\eq[(b'+a)][\Obj 0]$ داریم
$\eq[(b' + \Obj 0)][\Obj 0]$. بنا بر اصل موضوعهٔ~$!Q_4$ از~$\Th{Q}$،
$\eq[(b' + \Obj 0)][b']$ را داریم و ازاین‌رو $\eq[b'][\Obj 0]$.
اما بنا بر اصل موضوعهٔ~$!Q_2$، $\eq/[b'][\Obj 0]$ نیز داریم، که
تناقض است.

حالت ۲: برای یک~$c$، $\eq[a][c']$. اما آنگاه
$\eq[(b' + c')][\Obj 0]$ را داریم. بنا بر اصل موضوعهٔ~$!Q_5$،
$\eq[(b' + c)'][\Obj 0]$ را داریم که باز هم با اصل موضوعهٔ~$Q_2$
در تناقض است.
\end{proof}

\begin{lem}
\ollabel{lem:less-nsucc}
برای هر عدد طبیعی~$n$،
  \[
  \Th{Q} \Proves
  \lforall[x][(x < \num {n+1} \lif (\eq[x][\Obj 0] \lor \dots \lor
    \eq[x][\num n]))].
  \]
\end{lem}

\begin{proof}
از استقرا روی~$n$ استفاده می‌کنیم. نخست حالت پایه، یعنی $n = 0$، را
در نظر بگیریم. در این حالت باید نشان دهیم
$a < \num 1 \lif \eq[a][\Obj 0]$، برای~$a$ دلخواه. فرض کنید
$a < \num 1$. آنگاه بنا
بر اصل موضوع تعریف‌کنندهٔ~$<$، داریم
$\lexists[y][\eq[(y'+a)][\Obj 0']]$ (زیرا $\num 1 \ident \Obj 0'$).

فرض کنید $b$ دارای این خاصیت باشد؛ یعنی $\eq[(b'+a)][\Obj 0']$ را
داشته باشیم. باید $\eq[a][\Obj 0]$ را نشان دهیم. بنا بر اصل موضوعهٔ
$!Q_3$، یا $\eq[a][\Obj 0]$ را داریم یا~$c$ای هست که
$\eq[a][c']$. در حالت نخست چیزی برای نشان‌دادن نیست. پس فرض کنید
$\eq[a][c']$. آنگاه $\eq[(b' + c')][\Obj 0']$ را داریم. بنا بر اصل
موضوعهٔ~$!Q_5$ از~$\Th{Q}$، $\eq[(b'+c)'][\Obj 0']$ را داریم. بنا بر
اصل موضوعهٔ~$!Q_1$، $\eq[(b' + c)][\Obj 0]$ را داریم. اما این امر،
بنا بر اصل موضوعهٔ~$!Q_8$، به معنای $c < \Obj 0$ است و با
\olref{lem:less-zero} تناقض دارد.

اکنون به گام استقرا می‌پردازیم. حالت~$n+1$ را ثابت می‌کنیم و حالت~$n$
را فرض می‌کنیم. پس فرض کنید $a < \num {n+2}$. باز هم با استفاده از~$!Q_3$
می‌توانیم دو حالت را از هم متمایز کنیم: $\eq[a][\Obj 0]$ و اینکه برای
یک~$b$، $\eq[a][b']$. در حالت نخست،
$\eq[a][\Obj 0] \lor \dots \lor \eq[a][\num{n+1}]$ بی‌درنگ نتیجه
می‌شود. در حالت دوم، $b' < \num {n+2}$، یعنی
$b' < \num{n+1}'$ را داریم. بنا بر اصل موضوعهٔ~$!Q_8$، برای یک~$c$،
$\eq[(c'+b')][\num{n+1}']$. بنا بر اصل موضوعهٔ~$!Q_5$،
$\eq[(c'+b)'][\num{n+1}']$. بنا بر اصل موضوعهٔ~$!Q_1$،
$\eq[(c'+b)][\num{n+1}]$، و بنابراین
$b < \num{n+1}$ بنا بر اصل موضوعهٔ~$!Q_8$. بنا بر فرض استقرا،
$\eq[b][\Obj 0] \lor \dots \lor \eq[b][\num{n}]$. با منطق از این
$\eq[b'][\Obj 0'] \lor \dots \lor \eq[b'][\num{n}']$ را به دست
می‌آوریم، و بنابراین
$\eq[a][\num{1}] \lor \dots \lor \eq[a][\num{n+1}]$، زیرا
$\eq[a][b']$.
\end{proof}

\begin{lem}
  \ollabel{lem:trichotomy} برای هر عدد طبیعی~$m$،
  \[
  \Th{Q} \Proves
  \lforall[y][((y < \num{m} \lor \num{m} < y) \lor \eq[y][\num{m}])].
  \]
\end{lem}

\begin{proof}
با استقرا روی~$m$. نخست حالت $m=0$ را در نظر بگیرید.
$\Th{Q} \Proves
\lforall[y][(\eq[y][\Obj 0] \lor \lexists[z][\eq[y][z']])]$ بنا بر~$!Q_3$.
بگذارید $a$ دلخواه باشد. آنگاه یا $\eq[a][\Obj 0]$ یا برای یک~$b$،
$\eq[a][b']$. در حالت نخست،
$(a < \Obj 0 \lor \Obj 0 < a) \lor \eq[a][\Obj 0]$ را نیز داریم.
اما اگر $\eq[a][b']$، آنگاه $\eq[(b' + \Obj 0)][(a + \Obj 0)]$،
بنا بر منطق~$\eq$. بنا بر~$!Q_4$،
$\eq[(a + \Obj 0)][a]$؛ پس $\eq[(b' + \Obj 0)][a]$ و در نتیجه
$\lexists[z][\eq[(z' + \Obj 0)][a]]$ را داریم. بنا بر تعریف~$<$ در
$!Q_8$، $\Obj 0 < a$. اگر $\Obj 0 < a$، آنگاه
$(\Obj 0 < a \lor a < \Obj 0) \lor \eq[a][\Obj 0]$ را نیز داریم.

اکنون فرض کنید
\begin{align*}
  \Th{Q} & \Proves \lforall[y][((y < \num{m} \lor \num{m} < y) \lor
    \eq[y][\num{m}])]
  \intertext{و می‌خواهیم نشان دهیم}
  \Th{Q} & \Proves \lforall[y][((y < \num{m+1} \lor \num{m+1} < y) \lor
    \eq[y][\num{m+1}])].
\end{align*}
بگذارید $a$ دلخواه باشد. بنا بر~$!Q_3$، یا $\eq[a][\Obj 0]$ یا برای
یک~$b$، $\eq[a][b']$. در حالت نخست،
$\eq[\num{m}' + a][\num{m+1}]$ را بنا بر~$!Q_4$ داریم، و بنابراین
$a < \num{m+1}$ بنا بر~$!Q_8$.

اکنون حالت دوم، یعنی $\eq[a][b']$، را در نظر بگیرید. بنا بر فرض
استقرا، $(b < \num{m} \lor \num{m} < b) \lor
    \eq[b][\num{m}]$.

فصل نخست، $b < \num{m}$، (بنا بر~$!Q_8$) هم‌ارز است با
$\lexists[z][\eq[(z' + b)][\num{m}]]$. فرض کنید $c$ این خاصیت را
داشته باشد. اگر $\eq[(c' + b)][\num{m}]$، آنگاه
$\eq[(c' + b)'][\num{m}']$ نیز داریم. بنا بر~$!Q_5$،
$\eq[(c' + b)'][(c' + b')]$. ازاین‌رو $\eq[(c' +
b')][\num{m}']$. فرمول
$\lexists[u][\eq[(u' + b')][\num{m+1}]]$ را با تعمیم وجودی روی~$c$
به دست می‌آوریم، با توجه به اینکه
$\num{m}' \ident \num{m+1}$. بنابراین اگر
$b < \num{m}$، آنگاه $b' < \num{m+1}$ و پس $a < \num{m+1}$.

اکنون فرض کنید $\num{m} < b$؛ یعنی
$\lexists[z][\eq[(z' + \num{m})][b]]$. فرض کنید $c$ چنین~$z$ای باشد؛
یعنی $\eq[(c' + \num{m})][b]$. بنا بر منطق،
$\eq[(c' + \num{m})'][b']$. بنا بر~$!Q_5$،
$\eq[(c' + \num{m}')][b']$. چون $\eq[a][b']$ و
$\num{m}' \ident \num{m+1}$، داریم $\eq[(c' + \num{m+1})][a]$.
بنا بر~$!Q_8$، $\num{m+1} < a$.

در پایان، $\eq[b][\num{m}]$ را فرض کنید. آنگاه بنا بر منطق،
$\eq[b'][\num{m}']$ و بنابراین $\eq[a][\num{m+1}]$.

پس از هر فصل حالت مربوط به~$m$ و~$b$ می‌توانیم فصل متناظر برای~$m+1$
و~$a$ را به دست آوریم.
\end{proof}

\begin{prop}
\ollabel{prop:rep-minimization}
اگر $!A_g(x, z, y)$ تابع $g(x, z)$ را در~$\Th{Q}$ بازنمایی کند، آنگاه
\[
!A_f(z,y) \ident !A_g(y, z, \Obj 0) \land \lforall[w][(w < y \lif \lnot
  !A_g(w, z, \Obj 0))]
\]
تابع $f(z) = \umin{x}{[g(x, z) = 0]}$ را بازنمایی می‌کند.
\end{prop}

\begin{proof}
نخست نشان می‌دهیم اگر $f(n) = m$، آنگاه
$\Th{Q} \Proves !A_f(\num n, \num m)$؛ یعنی
\begin{align}
  \Th{Q} & \Proves !A_g(\num{m}, \num{n}, \Obj 0) \land \lforall[w][(w <
    \num{m} \lif \lnot !A_g(w, \num{n}, \Obj 0))]. \notag
  \intertext{چون $!A_g(x, z, y)$ تابع $g(x, z)$ را بازنمایی می‌کند و
    $g(m, n) = 0$ هرگاه $f(n) = m$، داریم}
\Th{Q} & \Proves !A_g(\num{m}, \num{n}, \Obj 0). \notag
\intertext{اگر $f(n) = m$، آنگاه برای هر $k < m$، $g(k, n) \neq 0$. پس}
\Th{Q} & \Proves \lnot !A_g(\num{k}, \num{n}, \Obj 0). \notag
\intertext{به دست می‌آوریم که}
\Th{Q} & \Proves \lforall[w][(w < \num{m} \lif \lnot
  !A_g(w, \num{n}, \Obj 0))]. \ollabel{rep-less}
\end{align}
اگر $m = 0$ باشد این امر بنا بر \olref{lem:less-zero} و در غیر این
صورت بنا بر \olref{lem:less-nsucc} نتیجه می‌شود.

اکنون نشان دهیم اگر $f(n) = m$، آنگاه $\Th{Q} \Proves
\lforall[y][(!A_f(\num{n}, y) \lif \eq[y][\num{m}])]$. باز هم
استدلال را به‌طور غیرصوری طرح می‌کنیم و صورت‌بندی صوری آن را به خواننده
می‌سپاریم.

فرض کنید $!A_f(\num{n}, b)$. از این (الف)~$!A_g(b, \num{n},
\Obj 0)$ و (ب)~$\lforall[w][(w < b \lif \lnot !A_g(w, \num{n}, \Obj
  0))]$ را به دست می‌آوریم. بنا بر \olref{lem:trichotomy}،
$(b < \num{m} \lor \num{m} < b) \lor \eq[b][\num{m}]$. نشان خواهیم
داد که هم $b < \num{m}$ و هم $\num{m} < b$ به تناقض می‌انجامند.

اگر $\num{m} < b$، آنگاه از~(ب) داریم
$\lnot !A_g(\num{m}, \num{n}, \Obj 0)$. اما $m = f(n)$؛ پس
$g(m, n) = 0$، و بنابراین
$\Th{Q} \Proves !A_g(\num{m}, \num{n}, \Obj 0)$، زیرا $!A_g$
تابع~$g$ را بازنمایی می‌کند. بنابراین به تناقض
می‌رسیم.

اکنون فرض کنید $b < \num{m}$. چون
$\Th{Q} \Proves \lforall[w][(w <
  \num{m} \lif \lnot !A_g(w, \num{n}, \Obj 0))]$ بنا بر
\olref{rep-less}، نتیجه می‌گیریم
$\lnot !A_g(b, \num{n}, \Obj 0)$. این نیز با~(الف) در تناقض است.
\end{proof}

\end{document}
